The Electron-Ion Collider (EIC) will further our understanding of the quarks and gluons which constitute visible matter by imaging the internal structure of protons and nuclei with greater precision than ever achieved. High-energy collisions at the EIC will lead to the production of heavy flavor particles such as $D^0$-mesons, which serve as probes to the gluon distributions within the colliding protons and neutrons. In order to develop efficient methods of analyzing these $D^0$-mesons, we explore the application of boosted decision trees (BDTs), a type of machine learning model, to separate $D^0$ signals from background based on features of the $D^0\rightarrow\pi+K$ decay topology. We present a preliminary analysis in which we apply an existing BDT framework to new simulated data, which demonstrates the negative impacts of machine-background effects on the $D^0$ analysis. We then optimize the BDT framework by limiting the kinematic range of the data, selecting more effective topological features, and tuning the hyperparameters which govern the BDT model structure. Finally, we demonstrate the advantage of BDT based analysis over the topological cuts determined through non-machine learning methods.